\documentclass[letterpaper, 12pt]{article}
% ============================================================
%  CONFIGURACION COMUN PARA SOLUCIONARIOS PEP 1
% ============================================================

\usepackage[utf8]{inputenc}
\usepackage[spanish]{babel}
\usepackage{amsmath, amssymb, amsfonts}
\usepackage{geometry}
\usepackage{fancyhdr}
\usepackage{enumitem}
\usepackage{graphicx}
\usepackage{hyperref}
\usepackage{needspace}
\usepackage{parskip}
\usepackage{xcolor}
\usepackage{tcolorbox}
\usepackage{eso-pic}
\usepackage{tikz}
\tcbuselibrary{skins, breakable}

\geometry{letterpaper, margin=1.7cm, top=3cm, bottom=2.5cm, headheight=2cm}

\definecolor{celesteSuave}{HTML}{F0F7FF}
\definecolor{azulFuerte}{HTML}{1A5276}
\definecolor{enlace}{HTML}{1A5276}

\newcommand{\R}{\mathbb{R}}
\newcommand{\Dom}{\operatorname{Dom}}
\newcommand{\Rec}{\operatorname{Rec}}
\newcommand{\Cod}{\operatorname{Cod}}

\newcommand{\logoup}{../../../../../template/images/logo_up.png}
\newcommand{\logousach}{../../../../../template/images/logo_usach.png}
\newcommand{\logofondo}{../../../../../template/images/logo_up.png}
\newcommand{\institucion}{Usach Premium}
\newcommand{\curso}{C\'alculo I}
\newcommand{\autor}{Jaime Riquelme}
\newcommand{\semestre}{1er. Semestre de 2026}
\newcommand{\fraseportada}{``Por y para estudiantes''}
\newcommand{\webinstitucion}{www.usachpremium.cl}
\newcommand{\instagraminstitucion}{https://www.instagram.com/usach.premium}
\newcommand{\transparencia}{0.05}

\hypersetup{
    colorlinks=true,
    linkcolor=enlace,
    urlcolor=enlace,
    citecolor=enlace,
    pdfborder={0 0 0},
}

\pagestyle{fancy}
\fancyhf{}
\fancyhead[L]{\includegraphics[height=1.2cm]{\logoup}}
\fancyhead[R]{%
    \begin{minipage}[b]{0.42\textwidth}
        \raggedleft
        \fontsize{9pt}{11pt}\selectfont
        \textbf{\color{azulFuerte}\institucion} \\
        \curso\ -- Solucionario PEP 1 \\
        \textit{\color{gray}@usach.premium}
    \end{minipage}\hspace{0.1cm}%
    \includegraphics[height=1.2cm]{\logousach}%
}
\fancyfoot[L]{\fontsize{9pt}{11pt}\selectfont \textit{\fraseportada}}
\fancyfoot[R]{\fontsize{9pt}{11pt}\selectfont P\'agina \thepage}
\renewcommand{\headrulewidth}{0.4pt}
\renewcommand{\footrulewidth}{0.4pt}

\fancypagestyle{plain}{
    \fancyhf{}
    \renewcommand{\headrulewidth}{0pt}
    \renewcommand{\footrulewidth}{0pt}
}

\newcommand{\MarcaAgua}{
    \AddToShipoutPicture{
        \AtPageCenter{
            \begin{tikzpicture}[remember picture, overlay]
                \node[opacity=\transparencia] at (0,0)
                    {\includegraphics[width=0.7\textwidth]{\logofondo}};
            \end{tikzpicture}
        }
    }
}

\newtcolorbox{solucionbox}[1]{
    colback=white,
    colframe=azulFuerte,
    coltitle=white,
    colbacktitle=azulFuerte,
    title=\textbf{#1},
    title after break=\textbf{#1} (continuaci\'on),
    fonttitle=\bfseries,
    arc=4pt,
    boxrule=1pt,
    enhanced,
    breakable,
    before={\par\Needspace{0.36\textheight}}
}

\newtcolorbox{respuestabox}{
    colback=celesteSuave,
    colframe=azulFuerte,
    boxrule=0.8pt,
    arc=4pt
}

\newcommand{\portadasolucionario}[2]{
    \thispagestyle{plain}
    \AddToShipoutPicture*{%
        \put(54, 700){\includegraphics[width=2cm]{\logoup}}
        \put(420, 706){\includegraphics[width=5cm]{\logousach}}
    }
    \vspace*{0.5cm}
    \begin{center}
        \color{azulFuerte}
        {\huge \textbf{SOLUCIONARIO}} \\[0.5cm]
        {\Huge \textbf{\curso}} \\[0.3cm]
        {\Large #1} \\[0.3cm]
        {\large #2} \\[1.4cm]
        \begin{tcolorbox}[colback=celesteSuave, colframe=azulFuerte, arc=10pt,
            width=0.82\textwidth, center, boxrule=1.5pt]
            \centering
            \Large\textbf{Desarrollo paso a paso}
        \end{tcolorbox}
    \end{center}
    \vspace{1.4cm}
    \begin{center}
        {\Large \textbf{\semestre}} \\[0.7cm]
        {\large \textbf{\autor\ -- \institucion}}
    \end{center}
    \vfill
    \begin{center}
        {\color{azulFuerte}\hrule height 0.5pt}
        \vspace{0.3cm}
        \begin{minipage}{0.45\textwidth}
            \centering
            \textbf{Instagram:} \\
            \small\href{\instagraminstitucion}{@usach.premium}
        \end{minipage}
        \begin{minipage}{0.45\textwidth}
            \centering
            \textbf{Web:} \\
            \small\href{http://\webinstitucion}{\webinstitucion}
        \end{minipage}
    \end{center}
    \newpage
    \MarcaAgua
}

\newcommand{\respuesta}[1]{
    \begin{respuestabox}
        \textbf{Respuesta:} #1
    \end{respuestabox}
}


\begin{document}
\portadasolucionario{Gu\'ia de Preparaci\'on PEP 1}{Secci\'on V -- Simulacros tipo PEP}

\section*{Simulacro A}

\begin{solucionbox}{Simulacro A -- Ejercicio 1}
Considere
\[
    f(x)=1+\sqrt{\frac{x-3}{x+2}}.
\]
\begin{enumerate}[label=\textbf{\alph*)}]
    \item Determine $\Dom(f)$.
    \item Determine $\Dom(f\circ f)$.
\end{enumerate}

\textbf{Soluci\'on:}

\textbf{a)} Para que $f$ est\'e definida, necesitamos
\[
    \frac{x-3}{x+2}\geq0,
    \qquad x\neq -2.
\]
Los puntos cr\'iticos son $x=-2$ y $x=3$.

\[
\begin{array}{c|ccc}
    & (-\infty,-2) & (-2,3) & (3,+\infty) \\ \hline
    x-3 & - & - & + \\
    x+2 & - & + & + \\ \hline
    \dfrac{x-3}{x+2} & + & - & +
\end{array}
\]

Por lo tanto,
\[
    \Dom(f)=(-\infty,-2)\cup[3,+\infty).
\]

\textbf{b)} Por definici\'on,
\[
    \Dom(f\circ f)=\{x\in\Dom(f):f(x)\in\Dom(f)\}.
\]
Como
\[
    f(x)=1+\sqrt{\frac{x-3}{x+2}},
\]
se tiene que
\[
    f(x)\geq 1
\]
para todo $x\in\Dom(f)$.

Entonces $f(x)$ no puede pertenecer al intervalo $(-\infty,-2)$. Por lo tanto, para que $f(x)\in\Dom(f)$, basta exigir
\[
    f(x)\in[3,+\infty).
\]
Es decir,
\[
    1+\sqrt{\frac{x-3}{x+2}}\geq3.
\]
Luego,
\[
    \sqrt{\frac{x-3}{x+2}}\geq2.
\]
Elevando al cuadrado:
\[
    \frac{x-3}{x+2}\geq4.
\]
Llevando todo a un lado:
\[
    \frac{x-3}{x+2}-4\geq0.
\]
Reducimos:
\[
    \frac{x-3-4(x+2)}{x+2}\geq0.
\]
Entonces
\[
    \frac{x-3-4x-8}{x+2}\geq0
    \quad\Longleftrightarrow\quad
    \frac{-3x-11}{x+2}\geq0.
\]
Multiplicando por $-1$:
\[
    \frac{3x+11}{x+2}\leq0.
\]
Los puntos cr\'iticos son
\[
    x=-\frac{11}{3},\qquad x=-2.
\]
La tabla de signos entrega
\[
    x\in\left[-\frac{11}{3},-2\right).
\]
Este intervalo ya est\'a contenido en $\Dom(f)$, por lo que
\[
    \Dom(f\circ f)=\left[-\frac{11}{3},-2\right).
\]

\respuesta{$\Dom(f)=(-\infty,-2)\cup[3,+\infty)$ y $\Dom(f\circ f)=\left[-\dfrac{11}{3},-2\right)$.}
\end{solucionbox}

\begin{solucionbox}{Simulacro A -- Ejercicio 2}
Considere $g:A\to B$ definida por
\[
    g(x)=-2+\sqrt{-x^2+8x-12}.
\]
\begin{enumerate}[label=\textbf{\alph*)}]
    \item Determine $A=\Dom(g)$.
    \item Determine un conjunto $A_1\subseteq A$ para que $g$ sea inyectiva y demu\'estrelo.
    \item Determine $B$ para que $g:A_1\to B$ sea biyectiva y encuentre $g^{-1}$.
\end{enumerate}

\textbf{Soluci\'on:}

\textbf{a)} Para que $g$ est\'e definida:
\[
    -x^2+8x-12\geq0.
\]
Multiplicamos por $-1$ y cambiamos el sentido:
\[
    x^2-8x+12\leq0.
\]
Factorizamos:
\[
    x^2-8x+12=(x-2)(x-6).
\]
Entonces
\[
    (x-2)(x-6)\leq0.
\]
La soluci\'on es
\[
    2\leq x\leq6.
\]
Por lo tanto,
\[
    A=\Dom(g)=[2,6].
\]

\textbf{b)} La expresi\'on bajo la ra\'iz puede escribirse como
\[
    -x^2+8x-12
    =4-(x-4)^2.
\]
Esto muestra que la funci\'on es sim\'etrica respecto de $x=4$. Por eso, para lograr inyectividad, escogemos una rama. Tomemos
\[
    A_1=[4,6].
\]

Probemos que $g$ es inyectiva en $A_1$. Sean $a,b\in[4,6]$ tales que
\[
    g(a)=g(b).
\]
Entonces
\[
    -2+\sqrt{-a^2+8a-12}
    =
    -2+\sqrt{-b^2+8b-12}.
\]
Sumando $2$:
\[
    \sqrt{-a^2+8a-12}
    =
    \sqrt{-b^2+8b-12}.
\]
Elevando al cuadrado:
\[
    -a^2+8a-12=-b^2+8b-12.
\]
Reduciendo:
\[
    -a^2+8a=-b^2+8b.
\]
Llevando todo a un lado:
\[
    b^2-a^2+8a-8b=0.
\]
Factorizamos:
\[
    (b-a)(b+a)-8(b-a)=0.
\]
Luego,
\[
    (b-a)(a+b-8)=0.
\]
Entonces
\[
    b-a=0
    \quad\vee\quad
    a+b-8=0.
\]
La primera opci\'on implica $a=b$. Para la segunda, como $a,b\in[4,6]$, se tiene $a+b\geq8$, y la igualdad $a+b=8$ ocurre solo cuando $a=b=4$. En ambos casos concluimos que
\[
    a=b.
\]
Por lo tanto, $g$ es inyectiva en $[4,6]$.

\textbf{c)} Para que $g:A_1\to B$ sea biyectiva, tomamos
\[
    B=\Rec(g|_{[4,6]}).
\]
Calculamos los extremos:
\[
    g(4)=-2+\sqrt{-16+32-12}=-2+\sqrt4=0,
\]
\[
    g(6)=-2+\sqrt{-36+48-12}=-2+\sqrt0=-2.
\]
Como en $[4,6]$ la funci\'on es decreciente,
\[
    \Rec(g|_{[4,6]})=[-2,0].
\]
Por lo tanto,
\[
    B=[-2,0].
\]

Ahora hallamos la inversa. Sea
\[
    y=-2+\sqrt{-x^2+8x-12}.
\]
Sumando $2$:
\[
    y+2=\sqrt{-x^2+8x-12}.
\]
Como $y\in[-2,0]$, se tiene $y+2\geq0$, y podemos elevar al cuadrado:
\[
    (y+2)^2=-x^2+8x-12.
\]
Usando la forma completada:
\[
    -x^2+8x-12=4-(x-4)^2.
\]
Entonces
\[
    (y+2)^2=4-(x-4)^2.
\]
Por lo tanto,
\[
    (x-4)^2=4-(y+2)^2.
\]
Como $x\in[4,6]$, se tiene $x-4\geq0$, luego
\[
    x-4=\sqrt{4-(y+2)^2}.
\]
Finalmente,
\[
    x=4+\sqrt{4-(y+2)^2}.
\]

\respuesta{$A=[2,6]$, una elecci\'on es $A_1=[4,6]$, $B=[-2,0]$ y $g^{-1}(y)=4+\sqrt{4-(y+2)^2}$.}
\end{solucionbox}

\begin{solucionbox}{Simulacro A -- Ejercicio 3}
Determine los valores de $k\in\R$ para que la par\'abola
\[
    y=(k-1)x^2+2x+k
\]
y la recta
\[
    y=kx-4
\]
se intersecten en dos puntos distintos.

\textbf{Soluci\'on:}

Igualamos la par\'abola con la recta:
\[
    (k-1)x^2+2x+k=kx-4.
\]
Llevamos todo a un lado:
\[
    (k-1)x^2+2x-kx+k+4=0.
\]
Reduciendo:
\[
    (k-1)x^2+(2-k)x+(k+4)=0.
\]

Para que haya dos intersecciones distintas, la ecuaci\'on debe ser cuadr\'atica y tener discriminante positivo.

Primero:
\[
    k-1\neq0
    \quad\Longrightarrow\quad
    k\neq1.
\]
Luego:
\[
    \Delta>0.
\]
Calculamos
\[
    \Delta=(2-k)^2-4(k-1)(k+4).
\]
Desarrollando:
\[
    (2-k)^2=k^2-4k+4,
\]
\[
    (k-1)(k+4)=k^2+3k-4.
\]
Entonces
\[
    \Delta=k^2-4k+4-4(k^2+3k-4).
\]
Luego,
\[
    \Delta=k^2-4k+4-4k^2-12k+16
    =-3k^2-16k+20.
\]
La condici\'on $\Delta>0$ queda:
\[
    -3k^2-16k+20>0.
\]
Multiplicando por $-1$:
\[
    3k^2+16k-20<0.
\]
Resolvemos la ecuaci\'on asociada:
\[
    3k^2+16k-20=0.
\]
Por f\'ormula cuadr\'atica:
\[
    k=\frac{-16\pm\sqrt{16^2-4(3)(-20)}}{2\cdot3}.
\]
Luego,
\[
    k=\frac{-16\pm\sqrt{256+240}}{6}
     =\frac{-16\pm\sqrt{496}}{6}.
\]
Como
\[
    \sqrt{496}=4\sqrt{31},
\]
tenemos
\[
    k=\frac{-16\pm4\sqrt{31}}{6}
     =\frac{-8\pm2\sqrt{31}}{3}.
\]
Como la par\'abola $3k^2+16k-20$ abre hacia arriba, la desigualdad
\[
    3k^2+16k-20<0
\]
se cumple entre sus ra\'ices:
\[
    \frac{-8-2\sqrt{31}}{3}<k<\frac{-8+2\sqrt{31}}{3}.
\]
Finalmente, excluimos $k=1$:
\[
    k\in\left(\frac{-8-2\sqrt{31}}{3},1\right)
    \cup
    \left(1,\frac{-8+2\sqrt{31}}{3}\right).
\]

\respuesta{$k\in\left(\dfrac{-8-2\sqrt{31}}{3},1\right)\cup\left(1,\dfrac{-8+2\sqrt{31}}{3}\right)$.}
\end{solucionbox}

\newpage
\section*{Simulacro B}

\begin{solucionbox}{Simulacro B -- Ejercicio 1}
Resuelva la siguiente inecuaci\'on:
\[
    \frac{3-|x+1|}{x^2-4}<0.
\]

\textbf{Soluci\'on:}

Primero factorizamos el denominador:
\[
    x^2-4=(x-2)(x+2).
\]
Luego, las restricciones son
\[
    x\neq -2,\qquad x\neq 2.
\]

Ahora encontramos los ceros del numerador:
\[
    3-|x+1|=0
    \quad\Longleftrightarrow\quad
    |x+1|=3.
\]
Entonces
\[
    x+1=3 \quad \vee \quad x+1=-3,
\]
por lo que
\[
    x=2 \quad \vee \quad x=-4.
\]

Los puntos cr\'iticos son
\[
    x=-4,\quad x=-2,\quad x=2.
\]
Observemos que $x=2$ anula numerador y denominador, por lo que no pertenece al dominio.

\[
\begin{array}{c|cccc}
    & (-\infty,-4) & (-4,-2) & (-2,2) & (2,+\infty) \\ \hline
    3-|x+1| & - & + & + & - \\
    x-2 & - & - & - & + \\
    x+2 & - & - & + & + \\ \hline
    \dfrac{3-|x+1|}{(x-2)(x+2)} & - & + & - & -
\end{array}
\]

Como buscamos que la expresi\'on sea menor que cero, tomamos los intervalos de signo negativo. No incluimos $x=-4$ porque la desigualdad es estricta, ni $x=-2,2$ porque anulan el denominador.

\respuesta{$x\in(-\infty,-4)\cup(-2,2)\cup(2,+\infty)$.}
\end{solucionbox}

\begin{solucionbox}{Simulacro B -- Ejercicio 2}
Considere
\[
    f(x)=\sqrt{\frac{5}{\sqrt{x+1}}-1},
    \qquad
    g(x)=x-4.
\]
\begin{enumerate}[label=\textbf{\alph*)}]
    \item Halle $\Dom(f)$.
    \item Halle $\Rec(f)$.
    \item Determine $f\circ g$ indicando expresamente su dominio y regla de asignaci\'on.
\end{enumerate}

\textbf{Soluci\'on:}

\textbf{a)} Para que $f$ est\'e definida, primero necesitamos que $\sqrt{x+1}$ est\'e en el denominador. Por lo tanto:
\[
    x+1>0
    \quad\Longleftrightarrow\quad
    x>-1.
\]
Adem\'as, el radicando exterior debe ser mayor o igual que cero:
\[
    \frac{5}{\sqrt{x+1}}-1\geq0.
\]
Entonces
\[
    \frac{5}{\sqrt{x+1}}\geq1.
\]
Como $\sqrt{x+1}>0$, multiplicamos:
\[
    5\geq \sqrt{x+1}.
\]
Elevando al cuadrado:
\[
    25\geq x+1
    \quad\Longleftrightarrow\quad
    x\leq24.
\]
Por lo tanto,
\[
    \Dom(f)=(-1,24].
\]

\textbf{b)} Si $x\in(-1,24]$, entonces
\[
    0<\sqrt{x+1}\leq5.
\]
Por lo tanto,
\[
    \frac{5}{\sqrt{x+1}}\geq1.
\]
Luego,
\[
    \frac{5}{\sqrt{x+1}}-1\geq0.
\]
Cuando $x=24$, se obtiene
\[
    f(24)=\sqrt{\frac{5}{5}-1}=0.
\]
Cuando $x\to -1^+$, se tiene $\sqrt{x+1}\to0^+$, por lo que la expresi\'on crece sin cota superior. Por lo tanto,
\[
    \Rec(f)=[0,+\infty).
\]

\textbf{c)} Por definici\'on:
\[
    \Dom(f\circ g)=\{x\in\Dom(g):g(x)\in\Dom(f)\}.
\]
Como $g(x)=x-4$ est\'a definida para todo real,
\[
    x-4\in(-1,24].
\]
Entonces
\[
    -1<x-4\leq24.
\]
Sumando $4$:
\[
    3<x\leq28.
\]
Por lo tanto,
\[
    \Dom(f\circ g)=(3,28].
\]
La regla de asignaci\'on es
\[
    (f\circ g)(x)=f(x-4).
\]
Luego,
\[
    (f\circ g)(x)=
    \sqrt{\frac{5}{\sqrt{(x-4)+1}}-1}
    =
    \sqrt{\frac{5}{\sqrt{x-3}}-1}.
\]

\respuesta{$\Dom(f)=(-1,24]$, $\Rec(f)=[0,+\infty)$, $\Dom(f\circ g)=(3,28]$ y $(f\circ g)(x)=\sqrt{\dfrac{5}{\sqrt{x-3}}-1}$.}
\end{solucionbox}

\begin{solucionbox}{Simulacro B -- Ejercicio 3}
Sea $k\in\R$ tal que
\[
    f:[-2,0]\to[k,12],
    \qquad
    f(x)=(x-1)^2+3
\]
est\'a bien definida.
\begin{enumerate}[label=\textbf{\alph*)}]
    \item Demuestre que $f$ es inyectiva.
    \item Determine $k$ para que $f$ sea sobreyectiva.
    \item Para ese valor de $k$, determine $f^{-1}$ indicando dominio, codominio y regla.
\end{enumerate}

\textbf{Soluci\'on:}

\textbf{a)} Sean $a,b\in[-2,0]$ tales que
\[
    f(a)=f(b).
\]
Entonces
\[
    (a-1)^2+3=(b-1)^2+3.
\]
Restando $3$:
\[
    (a-1)^2=(b-1)^2.
\]
Luego,
\[
    |a-1|=|b-1|.
\]
Como $a,b\in[-2,0]$, se tiene
\[
    a-1<0,\qquad b-1<0.
\]
Por lo tanto,
\[
    |a-1|=-(a-1),
    \qquad
    |b-1|=-(b-1).
\]
As\'i,
\[
    -(a-1)=-(b-1).
\]
Luego,
\[
    a-1=b-1,
\]
y por tanto
\[
    a=b.
\]
Entonces $f$ es inyectiva.

\textbf{b)} Calculamos el recorrido en $[-2,0]$:
\[
    f(-2)=(-2-1)^2+3=9+3=12,
\]
\[
    f(0)=(0-1)^2+3=1+3=4.
\]
Como la funci\'on es decreciente en $[-2,0]$,
\[
    \Rec(f)=[4,12].
\]
Para que $f:[-2,0]\to[k,12]$ sea sobreyectiva, debemos tener
\[
    [k,12]=[4,12].
\]
Por lo tanto,
\[
    k=4.
\]

\textbf{c)} Para $k=4$, la funci\'on es biyectiva. Hallamos su inversa:
\[
    y=(x-1)^2+3.
\]
Entonces
\[
    y-3=(x-1)^2.
\]
Como $x\in[-2,0]$, se tiene $x-1<0$, por lo que
\[
    \sqrt{y-3}=|x-1|=-(x-1)=1-x.
\]
Luego,
\[
    x=1-\sqrt{y-3}.
\]
Por lo tanto,
\[
    f^{-1}(y)=1-\sqrt{y-3}.
\]

\respuesta{$k=4$ y $f^{-1}(y)=1-\sqrt{y-3}$, con $f^{-1}:[4,12]\to[-2,0]$.}
\end{solucionbox}

\end{document}
